\clearpage
\begin{appendices}
\renewcommand{\thechapter}{\Asbuk{chapter}}
\titleformat{\chapter}[display]{\filcenter}{\MakeUppercase{Приложение} \thechapter}{0pt}{\bfseries}

\chapter{ИСХОДНЫЙ КОД ПРОЕКТА}

Полный исходный код вычислительного ядра \texttt{if97\_core}, сервиса и пользовательских приложений доступен в репозитории на GitHub.

\begin{figure}[htbp]
    \centering
    \qrcode[height=3cm]{https://github.com/veshchin/if97_calculator}
    \caption{QR-код для доступа к репозиторию проекта}
\end{figure}

\chapter{КОНТРОЛЬНЫЕ МАТЕРИАЛЫ И ФРАГМЕНТЫ КОДА}

\section{Контрольные точки и методика верификации}

Верификация вычислительного ядра выполнена как многоуровневый контур проверки, в котором эталонные данные и типы тестов разделены по назначению.  
Контрольный CSV-набор формировался из трёх классов источников: учебного пособия Гидаспова и Севериной~\cite{gidaspov_severina_2014}, справочника Александрова и Григорьева~\cite{alexandrov_grigorev_1999} и нормативных релизов IAPWS-IF97~\cite{iapws_if97_2012}, а для сложных обратных маршрутов региона~3 дополнительно использовались опубликованные backward equations для $v(p,T)$~\cite{iapws_vpt_region3}.  
Такое сочетание позволяет сопоставлять учебные табличные значения, независимую справочную сверку и формулы стандарта, не смешивая источники между собой.  
Актуальное состояние тестирования зафиксировано
в \path{TESTS.md}
(09.04.2026, commit \path{c76dbafa2e2f93f72d2ceae4eadd7f7de78ac4c3}).
Перегенерация отчёта выполняется скриптом \path{regenerate_tests_md.sh}.

Структура тестового контура разделена на функциональные, сервисные и нагрузочные уровни.
Модульные тесты проверяют отдельные формулы, вспомогательные функции и численные инварианты, например классификацию региона, поведение на границах областей и корректность обратных решателей.
Интеграционные тесты проверяют маршруты расчёта на уровне публичного фасада \texttt{If97} и сверяют результаты с эталонными значениями в допустимой погрешности.
Доктесты фиксируют работоспособность примеров публичного API и обеспечивают согласованность документации с кодом.
Отдельно выделены транспортные и сервисные тесты, которые поднимают реальный gRPC-сервер на локальном интерфейсе обратной петли, а также проверочные тесты серверной части Tauri без запуска графических диалогов.

\begin{table}[htbp]
    \centering
    \begin{tblr}{
        width=\linewidth,
        colspec={|p{3.1cm}|X|X|},
        row{1}={font=\bfseries},
        hline{1,2,Z},
    }
        Уровень теста & Что проверяется & Источник данных или инструмент \\
        Модульные тесты & Региональные формулы IF97, границы областей, инварианты и локальные вспомогательные функции & Нормативные релизы IAPWS-IF97, учебные и справочные контрольные точки \\
        Интеграционные CSV-тесты & Сквозные маршруты расчёта через \texttt{If97} и сравнение с эталонами & Единый CSV-набор, собранный из материалов IAPWS и табличных источников \\
        Доктесты & Работоспособность примеров публичного API и согласованность текста документации с кодом & Примеры из исходного кода и документации \\
        Smoke tests & Базовая жизнеспособность Tauri/native-контуров без запуска GUI & Headless-проверки стартовых сценариев \\
        Performance tests & Пропускная способность ядра и транспорта в debug/release-профилях & \texttt{load\_stress\_core\_csv}, \texttt{if97\_loadtest}, \texttt{grpc\_loadtest} \\
        Health/service tests & Доступность gRPC-сервиса, статусы \texttt{Serving/NotServing} и ограничения очередей & \texttt{tonic-health}, локальный gRPC-сервер и batch-валидация \\
    \end{tblr}
    \caption{Структура тестового контура проекта \texttt{if97\_calculator}}
    \label{tab:test_structure}
\end{table}

Таким образом, функциональная верификация опирается на контрольные значения и regression CSV, а нагрузочные сценарии и health-check-и отделены от неё организационно и по критерию прохождения.
Профиль \texttt{debug} используется для локальной диагностики и отладки, а \texttt{release} --- для финальной валидации и воспроизводимых измерений.

Отдельное внимание уделяется особым зонам области применимости.  
К таким зонам относятся линия насыщения (регион~4) и околокритическая область (регион~3), где свойства резко меняются, а обратные зависимости требуют устойчивых численных процедур.  
По состоянию контрольного прогона от 09.04.2026
для \texttt{if97\_core} подтверждены
18 модульных тестов, 1 интеграционный CSV-тест и 1 доктест.
Для \texttt{if97\_app\_api}
подтверждены тесты безопасной сериализации
\texttt{NaN}/\texttt{Inf} в JSON.
Для \path{if97_calculator_service}
подтверждены 5 модульных тестов,
включая проверки пакетной валидации, проверки работоспособности и ограничения давления очередей.
Для серверной части Tauri
подтверждены проверочные тесты
расчёта состояния, табличного парсинга
и генерации купола насыщения.
Все функциональные тесты завершились успешно.
	\texttt{\#[ignore]}-сценарии производительности
	запускаются отдельно и оформляются в отчёте как разделы замеров производительности.

\section{Пример использования API вычислительного ядра на Rust}

Ниже приведён укрупнённый пример вызовов публичного API ядра \texttt{if97\_core}.  
В реальном проекте в интерфейсных оболочках дополнительно выполняется разбор пользовательского ввода, валидация формата данных и отображение результатов.  

\begin{verbatim}
use if97_core::{If97, units::{MegaPascal, Kelvin}};

fn main() -> Result<(), if97_core::domain::If97Error> {
    let if97 = If97::new();

    // pT: direct calculation
    let state = if97.pt(
        MegaPascal(10.0),
        Kelvin(773.15),
    )?;
    println!(
        "h = {} kJ/kg, s = {} kJ/(kg*K)",
        state.h,
        state.s
    );

	    // ph: backward calculation
	    let state2 = if97.ph(
        MegaPascal(15.0),
        state.h,
    )?;
    println!(
        "T = {} K, region = {}",
        state2.t,
        state2.region
    );

    Ok(())
}
\end{verbatim}

\section{Фрагмент DTO для обмена между слоями}

Для взаимодействия между интерфейсной частью и серверным слоем используется отдельный крейт \texttt{if97\_app\_api}.  
Он содержит сериализуемые структуры, не зависящие от внутренних типов ядра, и обеспечивает стабильность контрактов обмена при развитии приложения.  
В частности, в DTO учитываются ограничения сериализации (например, необходимость корректного кодирования специальных числовых значений).  

\chapter{ДОПОЛНИТЕЛЬНЫЕ СХЕМЫ, ДИАГРАММЫ И МАНИФЕСТЫ}

\section{Схема потоков данных в программном комплексе}

В проекте реализован подход «одно ядро --- несколько потребителей», в котором \texttt{if97\_core} остаётся единственным источником вычислений.  
Укрупнённая схема потоков данных приведена ниже.  

\begin{verbatim}
                 +------------------+
                 |     if97_core     |
                 | (single compute)  |
                 +---------+---------+
                           |
         +-----------------+------------------+
         |                 |                  |
 +-------+------+   +------+--------+   +-----+------+
 |  FLTK UI     |   | Yew/WASM+Tauri|   | gRPC service|
 | (desktop)    |   |   (desktop)    |   | (batch API) |
 +--------------+   +---------------+   +-------------+
\end{verbatim}

\section{Фрагмент gRPC-протокола и SoA-представление данных}

Для сервисного контура применяется gRPC~\cite{grpc_docs}.  
Ниже приведён укрупнённый пример описания сообщений, использующих SoA-представление данных.  
Такой формат снижает накладные расходы при пакетной обработке и упрощает параллельные вычисления по данным.  

\begin{verbatim}
message PtBatchRequest {
  uint64 batch_id = 1;
  repeated double p = 2; // MPa
  repeated double t = 3; // K
}

message PtBatchResponse {
  uint64 batch_id = 1;
  repeated double h = 2; // kJ/kg
  repeated uint32 status = 3;
  uint32 batch_error_code = 4;
  string batch_error_message = 5;
}
\end{verbatim}

\section{Перечень параметров конфигурации сервиса}

Конфигурация сервисного слоя выполняется через переменные окружения, что упрощает развёртывание в контейнерной и оркестраторной среде.  
К основным параметрам относятся:
\begin{itemize}
\item параметр \path{IF97_GRPC_ADDR} --- адрес и порт gRPC-сервера;  
\item параметры \path{IF97_IO_THREADS} и \path{IF97_CPU_THREADS} --- настройка параллелизма ввода-вывода и вычислений;  
\item параметры \path{IF97_MAX_IN_FLIGHT_PER_CONN} и \path{IF97_MAX_IN_FLIGHT_GLOBAL} --- лимиты для ограничения давления очередей;  
\item параметры \path{IF97_MAX_BATCH_LEN} и \path{IF97_GRPC_MAX_MSG_BYTES} --- ограничения на размеры батчей и сообщений;  
\item параметр \path{IF97_DRAIN_SECS} --- время корректного завершения (штатного опустошения очереди) при остановке;  
\item параметр \path{IF97_KERNEL} --- выбор режима вычислительного ядра (\texttt{core} или \texttt{stub}).  
\end{itemize}

\section{Фрагменты инфраструктурных файлов для Docker, Kubernetes, непрерывной интеграции и доставки}

Контейнеризация сервиса выполнена с использованием Docker~\cite{docker_docs}.  
Эксплуатационная инфраструктура ориентирована на развёртывание в Kubernetes~\cite{kubernetes_docs} и на автоматизацию сборки и релизов через GitHub Actions~\cite{github_actions_docs}.  

Ниже приведены укрупнённые примеры фрагментов конфигураций, используемых в проекте.  

\textbf{Фрагмент ресурса развёртывания Kubernetes (Deployment):}
\begin{verbatim}
env:
  - name: IF97_GRPC_ADDR
    value: "0.0.0.0:50051"
  - name: IF97_CPU_THREADS
    value: "8"
\end{verbatim}

\textbf{Фрагмент конвейера GitHub Actions:}
\begin{verbatim}
on:
  push:
    tags:
      - "v*.*.*"
\end{verbatim}

\end{appendices}
